\section{Objetivos}\label{sec:objetivos}

\subsection{Objetivo general}
Desarrollar una red neuronal de clasificación binaria de alta letalidad en siniestros viales fatales del Perú, con pipeline completo (análisis exploratorio $\rightarrow$ preprocesamiento $\rightarrow$ ingeniería de características $\rightarrow$ entrenamiento $\rightarrow$ evaluación $\rightarrow$ interfaz gráfica), reproducible y documentado.

\subsection{Objetivos específicos}
\begin{itemize}
  \item \textbf{OE1.} Caracterizar la siniestralidad fatal nacional mediante análisis exploratorio de datos, con al menos ocho hallazgos interpretados.
  \item \textbf{OE2.} Construir al menos doce características derivadas de las variables pre-impacto del registro, cada una justificada.
  \item \textbf{OE3.} Entrenar un perceptrón multicapa (MLP) y compararlo contra tres líneas base: \texttt{DummyClassifier} (clase mayoritaria), regresión logística y \emph{random forest}. El MLP debe superar claramente al \texttt{DummyClassifier}; si no supera a la regresión logística o al \emph{random forest}, se reporta y discute honestamente.
  \item \textbf{OE4.} Identificar los cinco factores de mayor influencia en la alta letalidad mediante SHAP.
  \item \textbf{OE5.} Desplegar una interfaz gráfica en Streamlit que reutilice el pipeline de entrenamiento sin discrepancias.
\end{itemize}

\subsection{Métricas comprometidas antes de entrenar}
El dataset es desbalanceado: la clase multifatal es minoritaria (10.2\,\%). Por ello se fijaron las métricas antes de ver los datos, como hace un profesional:
\begin{itemize}
  \item Métrica primaria de selección de modelo: \textbf{F1-score de la clase multifatal}.
  \item Métricas de reporte obligatorio: \emph{precision} y \emph{recall} de la clase multifatal, PR-AUC, ROC-AUC, matriz de confusión completa y \emph{accuracy} solo como referencia.
  \item Se prohíbe argumentar calidad con \emph{accuracy} aislada: un modelo que responde siempre ``letalidad simple'' alcanza \emph{accuracy} alta y es inútil.
  \item Costo asimétrico: el falso negativo (subestimar un siniestro que dejó múltiples víctimas) es más costoso que el falso positivo, por lo que la calibración del umbral privilegia el \emph{recall}-multifatal.
\end{itemize}